\documentclass[11pt]{article}

\usepackage[utf8]{inputenc}
\usepackage[T1]{fontenc}

\usepackage[
	a4paper,
	left = 2cm,
	right = 2cm,
	bottom = 1cm, includefoot,
	top = 1.2cm, includehead,
]{geometry}

\usepackage{amsmath,amssymb,amsfonts}
\RequirePackage{physics}
\usepackage{siunitx}
\usepackage[version=4]{mhchem}
\usepackage{chemfig}
\AtBeginDocument{\RenewCommandCopy\qty\SI}
\NewDocumentCommand{\chem}{m}{\ce{#1}}

\usepackage{enumitem}
\usepackage{booktabs}

\numberwithin{equation}{section}
\numberwithin{figure}{section}

\usepackage{xcolor}
\usepackage[most]{tcolorbox}
\tcbuselibrary{breakable,skins}

\newtcolorbox{concept}[1][]{
	colback   = white,
	colframe  = gray!50,
	fonttitle = \small,
	enhanced,
	boxed title style = {
		colback  = white,
		colframe = white,
		boxrule  = 0pt,
		left     = 0.2em,
		right    = 0.1em,
	},
	breakable,
	boxrule = 0.5pt,
	left    = 6pt,
	right   = 6pt,
	sharp corners,
}

\newtcolorbox{phys-note}[1][]{
	colback   = white,
	breakable,
	boxrule = 0pt,
	left    = 6pt,
	right   = 6pt,
	top     = 0pt,
	bottom  = 0pt,
	sharp corners,
	enhanced, 
	borderline north = {0pt}{0pt}{white},
	borderline south = {0pt}{0pt}{white},
	borderline east  = {0pt}{0pt}{white},
	borderline west  = {1pt}{0pt}{black},
}

\newtcolorbox[auto counter, number within=section]{workedexample}[1][]{
	colback   = white,
	fonttitle = \bfseries,
	colframe  = white,
	title     = {Worked Example~\thetcbcounter#1},
	coltitle  = black,
	colbacktitle = white,
	attach boxed title to top left = {
		xshift = -3mm,
	},
	breakable,
	boxrule = 0pt,
	left    = 4mm,
	right   = 0pt,
	top     = 0pt,
	bottom  = 0pt,
	sharp corners,
	enhanced,
}

\usepackage{hyperref}

\title{\textbf{02 -- Chemical Language and Nomenclature}}
\author{Rodrigo Sánchez-Martínez}
\date{\today}

% --
\begin{document}
	
\maketitle	
	
%.
\tableofcontents

\section{Elements and Chemical Symbols}

	% --
	\subsection{The periodic table}

	\begin{concept}
		An \emph{element} is a substance whose atoms all share the same number of protons.
		The periodic table orders elements by increasing atomic number $Z$, arranged so that elements with similar valence configurations fall in the same column (group).
	\end{concept}

	\begin{phys-note}
		The recurring chemical behavior down a group is a direct consequence of the shell structure derived in Chapter~01:
		elements in the same group share the same valence electron configuration (e.g.\ $n\mathrm{s^1}$ for the alkali metals), so periodic trends in atomic radius, ionization energy, and electronegativity are simply different faces of the same underlying orbital filling pattern ($\mathrm{1s},\,\mathrm{2s},\,\mathrm{2p},\,\mathrm{3s},\,\dots$) dictated by the Aufbau principle and the Pauli exclusion principle.
	\end{phys-note}
	
	\subsection{Atomic number}

	\begin{concept}
		The \emph{atomic number} $Z$ is the number of protons in the nucleus.
		It uniquely identifies the element and equals the number of electrons in a neutral atom.
	\end{concept}

	\begin{phys-note}
		The electronic structure of an atom follows from the nuclear charge $+Ze$ acting on the electrons through the Coulomb potential
		\begin{equation}
			V(r) = -\frac{Ze^{2}}{4\pi\varepsilon_0 r},
		\end{equation}
		which enters the electronic Schrödinger equation discussed in Chapter~1.
		Changing $Z$ changes the element; changing the electron count while $Z$ is fixed only changes the ionization state.
	\end{phys-note}
	
	\subsection{Atomic mass}

	\begin{concept}
		The \emph{atomic mass unit} (u, or dalton, Da) is defined as $1/12$ the mass of a $^{12}\mathrm{C}$ atom:
		\begin{equation}
			\SI{1}{u} = \frac{1}{12}\, m(^{12}\mathrm{C}) \approx \SI{1.6605e-27}{kg}.
		\end{equation}
		The \emph{atomic mass} (or relative atomic mass, $A_r$) listed on the periodic table is the isotopic-abundance-weighted average
		\begin{equation}
			A_r = \sum_i f_i\, m_i,
		\end{equation}
		where $f_i$ is the natural fractional abundance of isotope $i$ and $m_i$ its isotopic mass in u.
	\end{concept}
	
	\begin{workedexample}[: average atomic mass of chlorine]
		Chlorine occurs naturally as \chem{^35Cl} ($m = \SI{34.969}{u}$, ${f = 0.7576}$) and \chem{^37Cl} ($m = \SI{36.966}{u}$, ${f = 0.2424}$). Then
		\begin{equation}
			A_r(\mathrm{Cl}) = (0.7576)(34.969) + (0.2424)(36.966) = \SI{35.45}{u},
		\end{equation}
		matching the tabulated value. Note that no single chlorine atom actually has this mass -- it is a statistical average over the isotopic population.
	\end{workedexample}
	
	\subsection{Chemical symbols}

	\begin{concept}
		Each element is denoted by a one- or two-letter symbol (e.g.\ H, He, Na, Fe), standardized by IUPAC.
		Many symbols derive from Latin or Greek names (Na from \textit{natrium}, Fe from \textit{ferrum}, W from \textit{wolfram}) rather than the modern English name.
	\end{concept}
	
	\subsection{Isotopic notation}

	\begin{concept}
		Isotopes of an element share the same $Z$ but differ in neutron number $N$, hence in mass number
		\begin{equation}
			A = Z + N.
		\end{equation}
		The standard notation is
		\begin{equation}
			\chem{^{$A$}_{$Z$}X}
		\end{equation}
		e.g.\ \chem{^2_1H} (deuterium) or \chem{^13_6C}.
	\end{concept}
	
	% --
	\section{Chemical Formulas}

	\subsection{Empirical formula}
	
	\begin{concept}
		The \emph{empirical formula} gives the simplest whole-number ratio of atoms of each element in a compound.
		It is obtained by dividing the moles of each element by the smallest mole value among them.
	\end{concept}

	\begin{workedexample}[: Empirical formula from mass composition]
		A compound is found to contain 40.0\% C, 6.7\% H, and 53.3\% O by mass. Taking \SI{100}{g} as basis,
		\begin{equation}
			n_{\mathrm C} = \frac{40.0}{12.01} = \SI{3.33}{mol}, \qquad
			n_{\mathrm H} = \frac{6.7}{1.008} = \SI{6.65}{mol}, \qquad
			n_{\mathrm O} = \frac{53.3}{16.00} = \SI{3.33}{mol}.
		\end{equation}
		Dividing by the smallest value ($3.33$) gives the ratio $\mathrm{C}:\mathrm{H}:\mathrm{O} = 1:2:1$, so the empirical formula is \chem{CH_2O}.
	\end{workedexample}

\end{document}